Estas son algunas preguntas que surgen de la práctica, con su respuesta.

\textbf{1. ¿Por qué el rayo no se dobla al pasar por la cara curva de la lente?}

Porque dentro de la lente viaja a lo largo de un radio del semicírculo, y un radio siempre llega perpendicular a la curva. Si el ángulo de incidencia es cero, la ecuación \eqref{eq:snell} da un ángulo de refracción también de cero.

\textbf{2. ¿Qué pasa con un ángulo de incidencia de 0 grados?}

El rayo atraviesa la lente sin desviarse en cualquiera de las dos orientaciones. Por eso se usó esa posición para alinear el equipo antes de empezar.

\textbf{3. ¿Por qué la reflexión total solo apareció con la luz entrando por la cara curva?}

Porque solo en ese caso la luz va de un medio lento, el acrílico, a uno rápido, el aire. Ahí el rayo se aleja de la normal, y llega un momento en que tendría que salir a más de 90 grados, lo cual es imposible, así que se refleja por completo. Al revés, del aire al acrílico, el rayo siempre se acerca a la normal y siempre puede entrar. Este fenómeno es el que permite que funcione la fibra óptica.

\textbf{4. ¿El índice de refracción depende del color de la luz?}

Sí, un poco. La luz violeta se dobla más que la roja en casi todos los materiales, y por eso un prisma separa la luz blanca en un arcoíris. Con una fuente de luz blanca como la de esta práctica, el rayo refractado puede verse ligeramente coloreado en los bordes, sobre todo cerca del ángulo crítico.
